
This appendix provides costed compatible joint realizations of the
splittable-pair constructions for all odd degrees except~$7$; for $n=7$ we
use a separate decodable septic base.  Algebraic splittability itself is not
exceptional at degree~$7$; the missing object is a cost-optimal joint
realization, as explained in \Cref{rem:canonical-split}.
The results hold over fields of characteristic zero and, for target degree
$n$, over fields of characteristic $p>n$.
The resulting degree-$n$ circuit uses exactly
$\lfloor n/2\rfloor+1$ multiplications and, after sharing affine forms,
at most
\[
   \min\!\left\{2n,\ \frac54n+6\lceil\log_2n\rceil^2+1\right\}
\]
additions or subtractions, with multiplicative height
$2\lceil\log_2 n\rceil+4$ for odd $n$ ($+5$ for even $n$, whose lift
adds one product) under the peeled scheduling of
\Cref{sec:peeled-Q}.  Thus the construction has leading arithmetic
cost $n/2$ multiplications and $5n/4$ additions on a critical path of
logarithmic length.  The preceding
decoder-calculus appendix supplies the common recovery language.  This appendix
certifies the construction, the multiplication count, and the logarithmic
height of the peeled scheduling; the separate addition-accounting appendix
proves the two stated addition bounds and the sequential baseline's height
from the same shared schedules.

\paragraph{Characteristic boundary.}
The recovery and circuit interfaces are characteristic-free.  The formulas
below are not: square peels and scalar pivots require~$2$ and several small
integers to be units.  A characteristic-$2$ treatment can reuse those
interfaces, but needs a separate recursive family rather than exceptional
branches of the $T_{k,2^l}$ recurrence.  Every helper that divides by~$2$
states that hypothesis explicitly.

\paragraph{Proof architecture.}
Although the formulas are recursive, the proof uses one decoder pattern
throughout:
\[
\begin{aligned}
  \text{construction identity}
    &\Longrightarrow \text{degree and support bounds}\\
    &\Longrightarrow \text{descending pivot table}
     \Longrightarrow \text{polynomial decoder}.
\end{aligned}
\]
The decoder calculus of \Cref{appendix:decoder-calculus} records polynomial
recovery and its causal version.  The fill construction supplies the known-powers gadgets
$Q_{2^s-1}$; the depth-balanced peeled variant of \Cref{sec:peeled-Q}
computes the same family at the identical ledger with height $s$.  The central $T_{k,2^l}$ lemma proves one
coefficient-triangular remainder invariant by three stage tables (even,
shared odd base, and ordinary odd).  The remaining lemmas attach small outer
gadgets and close a strong induction on the target degree.  Decoder summaries
after the longer proofs collect their steps; the proofs themselves establish
every pivot and information cutoff.

\paragraph{Construction atlas.}
The appendix can be read as the following small collection of circuit moves.
Each entry says both what the construction computes and how its inverse exposes
the hidden block.  On a first reading, this atlas together with the decoder
figures gives the logical spine; the intervening coefficient calculations
certify the stated arrows.
\begin{multicols}{2}
\small
\raggedcolumns
\noindent\textbf{Compatible closures.}
Sum, multiply, or square two pairs; descend through the translated windows
and peel the unique unresolved row.  These are the composition rules.

\smallskip\noindent\textbf{Fill $A_l$.}
Attach monic heads above and below a pair; decode upper head
$\to$ shifted input window $\to$ lower head.  Iterating this move gives the
Mersenne gadgets $Q_{2^s-1}$.

\smallskip\noindent\textbf{Peeled $Q^{\mathrm{peel}}$.}
Two half-size children under one keyed glue factor; decode by peeling
against the known monic $H$, reading $\gamma$ at the top row, and
subtracting.  Same ledger as the fill route, height $k$.

\smallskip\noindent\textbf{The $T_{k,2^l}$ recursion.}
Halve the exponent and attach a fresh tail; decode tail pivots
$\to$ shifted recursive block $\to$ low $Q$-block.  This is the central
compatible-pair construction.

\smallskip\noindent\textbf{The $Q_{4k+1}$ crown.}
Place a $T_{2k,2}$ pair under an affine crown; read five crown pivots and
then run the recovered $T$ decoder.  This handles degrees $1\pmod4$.

\smallskip\noindent\textbf{The barred gadget $\bar Q_{8k+7}$.}
Place a $T_{k,8}$ pair inside one level-$4$ fill; decode its top scalar/block
pivots $\to$ two seam rows $\to$ inner $T$ block $\to$ low rows.

\smallskip\noindent\textbf{Outer induction.}
Wrap a smaller pair in one or two difference-of-squares shells; peel the
shells, decode the smaller pair and its recorded powers, and only then decode
the auxiliary gadgets.

\smallskip\noindent\textbf{Finite bases.}
The fused circuits at $7,15,27,31$ use the same top-down shell and pivot
rules; they bridge the degrees at which a generic branch would miss the
optimal multiplication count.
\end{multicols}

\begin{figure}[H]
  \centering
  \resizebox{\textwidth}{!}{% The visual grammar used throughout the construction appendix.
\begin{tikzpicture}[x=1cm,y=1cm,>=Latex]
  % Panel (a): one causal pivot.
  \node[fp panel title] at (0,4.05) {(a) One causal pivot};
  \draw[fp flow] (0,3.62) -- (6.25,3.62);
  \node[fp label, anchor=east] at (0,3.62) {high degree};
  \node[fp label, anchor=west] at (6.25,3.62) {low degree};

  \node[fp known, minimum width=28mm] (seen) at (1.4,2.95)
    {higher observed rows};
  \node[fp active, minimum width=14mm, right=0pt of seen] (row)
    {$[x^j]\Phi$};
  \node[fp unread, minimum width=22mm, right=0pt of row] (unread)
    {rows $<j$};

  \node[fp known, minimum width=28mm] (highpar) at (1.4,1.82)
    {$\alpha_{>j}$ already decoded};
  \node[fp pivot, minimum width=14mm, right=0pt of highpar] (par)
    {$\alpha_j$};
  \node[fp unread, minimum width=22mm, right=0pt of par]
    {$\alpha_{<j}$};
  \draw[fp decode] (row.south) -- node[fp label, right, text=fpBlue!80!black]
    {subtract $F_j$; divide by $\lambda_j$} (par.north);
  \node[fp label, anchor=west] at (0,0.92)
    {$[x^j]\Phi=\lambda_j\alpha_j+F_j(\alpha_{>j})$,
     \quad$\lambda_j\ne0$};

  % Panel (b): nested shells.
  \node[fp panel title] at (7.45,4.05) {(b) A top-down construction};
  \draw[fp decode] (7.45,3.62) -- node[fp label, above]
    {peel, subtract, expose} (13.45,3.62);

  \node[fp active, minimum width=31mm, anchor=west] (outer) at (7.45,2.95)
    {outer monic shell};
  \node[fp seam, minimum width=7mm, anchor=west] at (10.55,2.95)
    {$\pm1$};
  \node[fp unread, minimum width=21mm, anchor=west] at (11.25,2.95)
    {hidden tail};

  \node[fp second, minimum width=27mm, anchor=west] (inner) at (9.35,1.82)
    {next shell};
  \node[fp seam, minimum width=7mm, anchor=west] at (12.05,1.82)
    {seam};

  \node[fp recursive, minimum width=25mm, anchor=west] (rec) at (11.05,0.92)
    {smaller instance};
  \draw[fp dependence] (outer.south east) to[bend left=15] (inner.north west);
  \draw[fp dependence] (inner.south east) to[bend left=15] (rec.north west);

  % Complete key.  Labels make the figure readable without colour.
  \node[fp active, minimum width=8mm, minimum height=4mm] at (0.4,0.2) {};
  \node[fp label, anchor=west] at (0.9,0.2) {being recovered};
  \node[fp second, minimum width=8mm, minimum height=4mm] at (3.35,0.2) {};
  \node[fp label, anchor=west] at (3.85,0.2) {next block};
  \node[fp recursive, minimum width=8mm, minimum height=4mm] at (5.85,0.2) {};
  \node[fp label, anchor=west] at (6.35,0.2) {smaller instance};
  \node[fp pivot, minimum width=8mm, minimum height=4mm] at (9.1,0.2) {};
  \node[fp label, anchor=west] at (9.6,0.2) {constant-slope pivot};
  \node[fp known, minimum width=8mm, minimum height=4mm] at (0.4,-0.45) {};
  \node[fp label, anchor=west] at (0.9,-0.45) {known};
  \node[fp unread, minimum width=8mm, minimum height=4mm] at (3.35,-0.45) {};
  \node[fp label, anchor=west] at (3.85,-0.45) {not yet read};
  \node[fp seam, minimum width=8mm, minimum height=4mm] at (5.85,-0.45) {};
  \node[fp label, anchor=west] at (6.35,-0.45) {boundary seam};
\end{tikzpicture}
}
  \caption{The visual language used in this appendix.  Coefficient degree
  decreases from left to right, which is also the order in which a decoder
  runs.  A causal pivot reads the current observed row and quantities decoded
  strictly earlier, but never a lower row.  Larger constructions are drawn as
  nested monic shells: peel the highest shell, subtract it (including any
  displayed boundary correction), and continue with the newly exposed
  block.  Colours distinguish blocks, while labels and line styles carry the
  same information in monochrome.}
  \label{fig:decoder-language}
\end{figure}

For concision, call a field $d$-\emph{admissible} if it has characteristic
zero or characteristic $p>d$.  This is a convenient sufficient condition,
not the exact one (\Cref{rem:char2}).  Every decoder displayed in this
appendix is a composition of unit pivots, monic divisions, and affine
pivots, root extractions and block solves whose slopes, root indices and
determinants are fixed integers; it is therefore defined, and inverts the
construction, over every field in which those integers are units.
Replaying the routing of \Cref{alg:final-construction} lists them exactly.
Along the recursive calls made for degree~$n$, every integer divided by is
one of
\begin{equation}
\label{eq:bad-primes}
  2,\qquad 2k,\qquad k\ (k\text{ even})\ \text{or}\ k(k-1)\ (k\text{ odd}),
  \qquad k^2,
\end{equation}
for an internal index $k$, and these occur as follows.
\begin{itemize}
  \item $2$: every monic square root, square gadget and scalar shift
    (\Cref{lem:monic-from-power,lem:square-gadget-boundary,lem:scalar-shift-square}),
    the septic base (\Cref{lem:septic-base}) and the finite bases $15,27,31$
    (\Cref{lem:special-cases-splittable}).
  \item $2k$: the five crown slopes $2k,2k,-2k,k,k$ of the $4k+1$ branch
    (\Cref{lem:4k+1-splittable}) and of $Q_{4k+1}(x,H_2)$
    (\Cref{lem:Q4k+1-from-H2}), and the slope $M=2k$ of a perturbed
    $T_{2k,2^l}$ call (\Cref{lem:causal-perturbed-T}).
  \item $k$ or $k(k-1)$: the $T$ recursion at index $k$
    (\Cref{lem:Rk2l,lem:Rk2l-leading-coeff}), followed down its halving
    chain $k\mapsto\lfloor k/2\rfloor$.  For even $k$ the stage table has
    slopes $-k$ and $k/2$ together with one square root and one scalar
    shift; for odd $k$ it has slopes $-k(k-1)$, $-(k-1)$ and $(k-1)/2$; the
    leading coefficient $\gamma_k$ divides $k(k-1)$ in both cases.
  \item $k^2$: the block determinant $\det M=-k^2$ and the two slopes $k$
    of the barred gadget $\bar Q_{8k+7}$ (\Cref{lem:barQ8k+7}).
\end{itemize}
The seam corrections $\gamma_m$,
$\binom m2$, $\tau=k(k-1)(k-2)/3=2\binom k3$ and $\theta_m$ are integers
that are subtracted, never divided by, and the fill gadgets $A_{2^l}$, the
known-powers gadgets $Q_{2^t-1}$ with their peeled variant, the gadgets
$Q_1,Q_3,Q_7$, the degree-$3$ base and the even lift use only unit pivots
and monic division.  Let $\mathrm{Bad}(n)$ be the set of primes dividing at
least one integer in \eqref{eq:bad-primes} along the routing for degree
$n$.  Then the construction and the decoder of $P_n$ are valid over every
field in which the elements of $\mathrm{Bad}(n)$ are units, i.e.\ whose
characteristic is not in $\mathrm{Bad}(n)$.  Every internal index $k$ in
\eqref{eq:bad-primes} is at most $(n-1)/4$, so every odd element of
$\mathrm{Bad}(n)$ is at most $(n-1)/4$ and $n$-admissibility is sufficient.
The script \texttt{tools/bad\_primes.py} replays the routing and prints
the pivot list; for instance
\[
\begin{aligned}
  &\mathrm{Bad}(5)=\mathrm{Bad}(9)=\mathrm{Bad}(15)=\mathrm{Bad}(17)
    =\mathrm{Bad}(31)=\{2\},\\
  &\mathrm{Bad}(13)=\mathrm{Bad}(27)=\{2,3\},\qquad
  \mathrm{Bad}(21)=\{2,5\},\qquad
  \mathrm{Bad}(29)=\{2,3,7\}.
\end{aligned}
\]
Thus $\mathrm{Bad}(n)=\{2\}$ for $n\in\{5,\dots,12\}$ and, e.g., for
$n=31,\,47,\,79$.  As a check, the reference decoder
\texttt{tools/polychain.py} completes its encode--decode round trip over
$\mathrm{GF}(p)$ for every odd $n\le75$ (and selected $n\le127$) and
every prime $5\le p\le31$ with $p\notin\mathrm{Bad}(n)$, including many
pairs with $p<n$; for $p\in\mathrm{Bad}(n)$ the displayed decoder divides
by zero.  (At $p=3$ that implementation additionally evaluates $\tau$ by a
division, which the paper's decoder does not need.)
No characteristic restriction beyond \eqref{eq:bad-primes} is hidden in
the recursion.
